% id: 08
% book: 2026 정의여고 1-2 중간고사 프린트 · 창의 변형 1세트 · 01 평면좌표·직선·원·도형의 이동
% section: p3
% figure: none
% answer: ④
% answer_source: 계산
% variant_level: 2 (원본 전사)
% 이 파일은 build.mjs 가 items.json 에서 만든다. 고칠 때는 items.json 을 고친다.
\smash{\raisebox{3.5mm}{\dmkichul{2026 정의여고 1-2 중간고사 프린트 · 창의 변형 1세트 08번 · 3쪽}}}
점~$(1{,}\mkern3mu 2)$를 지나고 기울기가 $k$인 직선~$l$이 있다. 점~$(4{,}\mkern3mu 3)$과 직선~$l$ 사이의 거리가 $\sqrt5$일 때, 가능한 모든 $k$의 값의 합은?
\begin{choices32}{\rule[-3ex]{0pt}{8ex}$-\dfrac12$}{\rule[-3ex]{0pt}{8ex}$\dfrac12$}{\rule[-3ex]{0pt}{8ex}$1$}{\rule[-3ex]{0pt}{8ex}$\dfrac32$}{\rule[-3ex]{0pt}{8ex}$2$}\end{choices32}
